\begin{spacing}{1.0} 
\begin{table} \centering \caption{Effects of Credit-Eligible Videos on Grades} 
\label{tab:secondstage_credit} 
\begin{threeparttable} 
\begin{tabular}{m{0.35\linewidth} *{4}{>{\centering\arraybackslash}m{0.1\linewidth}}}
\toprule
 & (1) & (2) & (3) & (4) \\
\midrule
\multicolumn{5}{l}{\textbf{Panel A}: Midterm 2 score} \\ 
\customlinespace \indentrow{RF: Incentive} & 0.176\sym{*} & 0.183\sym{*} & 0.176\sym{*} & 0.174\sym{*} \\
 & (0.090) & (0.094) & (0.090) & (0.096) \\
\customlinespace \indentrow{2SLS: 10 videos} & 0.320\sym{*} & 0.329\sym{*} & 0.366\sym{**} & 0.350\sym{*} \\
 & (0.176) & (0.183) & (0.181) & (0.191) \\
\customlinespace \indentrow{2SLS: 1 hour of videos} & 0.197\sym{*} & 0.201\sym{*} & 0.225\sym{**} & 0.210\sym{*} \\
 & (0.108) & (0.111) & (0.111) & (0.115) \\
\midrule 
Observations & 395 & 362 & 395 & 362 \\
\midrule 
\multicolumn{5}{l}{\textbf{Panel B}: Final exam score} \\ 
\customlinespace \indentrow{RF: Incentive} & 0.175\sym{**} & 0.174\sym{*} & 0.175\sym{**} & 0.138 \\
 & (0.089) & (0.103) & (0.088) & (0.103) \\
\customlinespace \indentrow{2SLS: 10 videos} & 0.093\sym{**} & 0.091\sym{*} & 0.094\sym{**} & 0.076 \\
 & (0.047) & (0.054) & (0.047) & (0.055) \\
\customlinespace \indentrow{2SLS: 1 hour of videos} & 0.059\sym{**} & 0.058\sym{*} & 0.060\sym{**} & 0.048 \\
 & (0.030) & (0.034) & (0.030) & (0.035) \\
\midrule 
Observations & 374 & 332 & 374 & 332 \\
Treatment assignment controls & Yes & No & Yes & Yes \\
Demographic controls & No & No & Yes & Yes \\
Pair Fixed Effects & No & No & No & Yes \\
\bottomrule
\end{tabular}
\Fignote{This table replicates Table~\ref{tab:secondstage_table} using credit-eligible video measures as the endogenous variable. The reduced-form (\textit{RF}) estimates do not depend on the video measure and are identical to Table~\ref{tab:secondstage_table}. Video measures count only credit-eligible video openings: a video counts only if the student did not open another video before this one could have been watched in its entirety (allowing one minute for timestamp rounding), the condition \textit{Incentive} students were told was required for credit. The rule is applied identically to all arms. Test scores are measured in standard deviation units. Model (1) contains linear controls for midterm 1 score and year; (2) is the difference in means and standard errors calculated using the repeated sampling framework of Neyman (1923); (3) and (4) use the post-double-selection (PDS) procedure of \textcite{bch2014a} to select control variables then estimate treatment effects and standard errors. Models (2) and (4) include only students whose matched-pair did not attrite from the experiment. \Regnote}
\end{threeparttable}
\end{table} 
\end{spacing}
